% !TEX root = Frozen Black Beans.tex
\documentclass[14pt]{extarticle}
\usepackage[utf8]{inputenc}
% --- REMOVE PAGE NUMBERING ---
\pagestyle{empty}
% --- MARGIN SETUP ---
\usepackage[left=0.75in, right=0.75in, top=0.5in, bottom=0.5in]{geometry}
% --- LIST SETTINGS ---
\usepackage{enumitem}
\setlist[itemize]{leftmargin=1em}
\setlist[itemize]{leftmargin=1em, itemsep=0.2em, label={\raisebox{-0.1ex}{\large\textbullet}}}
\setlist[enumerate]{leftmargin=1.5em, label=\textbf{\arabic*.}, itemsep=0.25em}
\usepackage{parskip}
% --- CONTROL SECTION FONT SIZE ---
\usepackage{titlesec}
\titleformat{\section}{\large\bfseries}{}{0em}{}
\titleformat{\subsection}{\normalsize\bfseries}{}{0em}{}
\titlespacing*{\subsection}{0pt}{0.75em}{0.25em}
\titlespacing*{\section}{0pt}{1em}{0.5em}
% --- FONT SETUP ---
\usepackage{helvet}
\renewcommand{\familydefault}{\sfdefault}
% --- CUSTOM COMMAND FOR TITLE ---
\newcommand{\makeRecipeTitle}[1]{
    {\noindent \Large \textbf{#1}}
    \par \vspace{-0.75em}
}
% --- CUSTOM COMMAND FOR SERVINGS ---
% Usage: \recipeServings{4--5} or \recipeServings{6 (one burrito each)}
% Place directly after the first \makeRecipeTitle on page 1.
\newcommand{\recipeServings}[1]{
    {\noindent \itshape Serves #1 \par}
    \vspace{-0.75em}
}
% --- CUSTOM COMMAND FOR SOURCE CREDIT ---
% Usage: \recipeSource{Creator -- Video Title}{https://...}
% Place just before \end{document}; pins the credit to the bottom of the last page.
\usepackage[hidelinks]{hyperref}
\urlstyle{same}
\newcommand{\recipeSource}[2]{
    \vfill
    {\noindent \footnotesize \textbf{Source:} #1 \\ \url{#2}}
}


\begin{document}

% ================= PAGE 1 =================
\makeRecipeTitle{Frozen Black Beans}
\recipeServings{8--10 (as a side)}

\section*{Ingredients}

\begin{itemize}
      \item Dried Black Beans – 12 oz
      \item Yellow Onion – 1, peeled and halved
      \item Garlic Cloves – 2, peeled
      \item Ginger – 1 knob, peeled
      \item Jalapeño Pepper – 1, deseeded
      \item Water – enough to cover the beans
      \item Salt and Pepper – to taste
\end{itemize}

\section*{Notes}

\begin{itemize}
      \item \textbf{Seasoning:} Hold the salt until the beans are fully cooked. Seasoning early makes the finished salt level much harder to judge.
      \item \textbf{Texture:} Blend before freezing for a refried-style base, or freeze the beans whole to keep your options open later.
\end{itemize}

\newpage
\makeRecipeTitle{Frozen Black Beans}
\section*{Instructions}

\begin{enumerate}

      \item \textbf{Load the Pressure Cooker} \\
            Add the beans, onion, garlic, ginger, and jalapeño to a pressure cooker and cover completely with water.

      \item \textbf{Pressure Cook} \\
            Cook for 18 minutes, or until the beans are fork tender. Leave them unseasoned at this stage.

      \item \textbf{Release the Steam} \\
            Once the beans are fully cooked, release the steam from the pressure cooker before opening the lid.

      \item \textbf{Season and Blend} \\
            Add salt and pepper to taste, then blend lightly with an immersion blender, leaving plenty of whole beans behind.

      \item \textbf{Adjust Thickness} \\
            If the beans are thinner than you like, return the pot to the burner and reduce until they reach your preferred thickness.

      \item \textbf{Portion and Freeze} \\
            Ladle into silicone molds and freeze overnight.

\end{enumerate}

\end{document}
